\section{Sviluppi futuri}
\label{sviluppi_futuri}

In questo capitolo vengono illustrate alcune possibili direzioni di sviluppo
della piattaforma, che si collocano su due piani complementari: da un lato la
risoluzione di alcuni problemi residui emersi durante l'esperimento con gli
utenti e dalla valutazione euristica, dall'altro l'estensione del sistema con nuove funzionalità
volte ad ampliarne le potenzialità e ad arricchire l'esperienza di studio. Per
ciascuna proposta vengono esposte la motivazione che la giustifica e i benefici
attesi in termini di usabilità complessiva.

\begin{itemize}
  \item \textbf{Personalizzazione e ottimizzazione dell'esperienza d'uso.}
  L'analisi dei dati ha confermato l'efficacia della riprogettazione; tuttavia,
  negli scenari più articolati (S1 e S2) alcuni utenti hanno descritto il flusso
  come ancora «macchinoso» per i ripetuti spostamenti tra sezioni, segnalando
  inoltre l'assenza di un supporto contestuale all'orientamento. A queste
  osservazioni si aggiungono due limiti rilevati in fase di valutazione
  euristica e legati al controllo dell'utente sui contenuti: gli output generati
  dagli strumenti di intelligenza artificiale non sono modificabili (P68) e i
  testi inseriti dagli utenti non supportano alcuna formattazione di base (P71).
  Si propongono pertanto quattro interventi complementari:
  \begin{itemize}
    \item \textbf{Azioni dirette e contestuali}: per ridurre i passaggi nei
    flussi più frequenti, una scorciatoia «Aggiungi
    a esame» dai salvataggi personali e dalla community (suggerita nello
    scenario S1).
    \item \textbf{Assistente conversazionale basato su IA}: un chatbot/AI
    assistant integrato nelle varie sezioni, capace di rispondere in linguaggio
    naturale e di guidare l'utente durante l'esecuzione dei task, a beneficio in
    particolare di chi incontra difficoltà di orientamento.
    \item \textbf{Modifica degli output generati dall'IA}: i contenuti prodotti
    dagli strumenti AI , oggi presentati come output
    «finali», diventerebbero modificabili e salvabili: l'utente potrebbe così
    correggere eventuali errori di trascrizione e integrare il testo con le
    proprie note, riacquisendo il controllo su contenuti che oggi è costretto
    ad accettare così come sono.
    \item \textbf{Formattazione di base dei testi}: le descrizioni degli appunti
    e i commenti verrebbero arricchiti con strumenti di formattazione inline e
    scorciatoie da tastiera per grassetto, corsivo ed elenchi, permettendo di
    strutturare i contenuti più lunghi e di evidenziare i passaggi rilevanti, a
    beneficio della leggibilità complessiva.
  \end{itemize}
  Nel loro insieme questi interventi mirano a snellire i percorsi di interazione
  più ricorrenti e ad ampliare il controllo dell'utente sui contenuti, riducendo
  il numero di click necessari e il carico cognitivo legato all'uso
  dell'interfaccia.

  \item \textbf{Caricamento ibrido e organizzazione dei documenti.}
  L'esperimento con gli utenti e la valutazione euristica hanno fatto emergere
  due criticità complementari legate alla gestione dei documenti: la rigidità
  dell'attuale flusso di caricamento, in cui la creazione di una cartella e
  l'inserimento di un file costituiscono due passaggi distinti e sequenziali, e
  l'assenza di strumenti organizzativi nella sezione Community, segnalata già in
  fase di valutazione euristica (P72). Si propongono pertanto due interventi:
  \begin{itemize}
    \item \textbf{Flusso di caricamento ibrido}: per collocare un documento in
    una nuova cartella l'utente deve oggi prima crearla e solo successivamente
    caricarvi il file, un vincolo percepito come rigido; l'esigenza di unificare
    le due operazioni è emersa in modo indipendente in due diverse attività
    dell'esperimento, configurandosi come la richiesta di miglioramento più
    ricorrente. All'interno della finestra di caricamento, un campo
    «Destinazione» con un'opzione «+ Nuova cartella» consentirebbe di creare la
    cartella e di caricarvi contestualmente il file, senza interrompere il
    flusso ed eliminando un passaggio superfluo in un'operazione frequente.
    \item \textbf{Sistema di gestione documentale per la Community}: la sezione
    presenta i contenuti condivisi come un
    elenco piatto, privo di cartelle e sottocartelle personalizzate, una
    struttura che con l'aumento dei materiali rischia di rendere la
    consultazione ingestibile nel lungo periodo. Coerentemente con la scelta
    progettuale di non replicare nella Community le cartelle personali della
    sezione Appunti, si propone l'introduzione di un sistema di gestione
    documentale condiviso, che organizzi gli appunti della community in raccolte
    gerarchiche per corso di studi, insegnamento e argomento, affiancate da tag
    e metadati, mantenendo la ricerca e la navigazione efficienti anche al
    crescere dei contenuti.
  \end{itemize}
  I due interventi agiscono sui due estremi del ciclo di vita del documento --
  il caricamento e la consultazione -- riducendo la rigidità segnalata dagli
  utenti e garantendo la scalabilità organizzativa della piattaforma nel lungo
  periodo.

  \item \textbf{Ridisegno del flusso di rinomina.} L'analisi dei dati ha
  evidenziato un'unica regressione rispetto alla versione originale, relativa
  alla rinomina degli elementi. Dai commenti degli utenti emergono tre cause
  precise: l'assenza di un'azione esplicita di rinomina, nascosta all'interno di
  un generico comando «Modifica»; l'apertura del campo già vuoto anziché
  precompilato, un evidente \textit{gulf of execution} in cui l'utente è
  costretto a ricostruire da sé un'informazione che il sistema già possiede; e
  l'impossibilità di rinominare direttamente dall'anteprima del documento. Si
  propone pertanto un ridisegno del flusso, articolato in tre interventi
  corrispondenti:
  \begin{itemize}
    \item \textbf{Azione di rinomina dedicata}: una voce «Rinomina» esplicita nel
    menu contestuale, distinta dal generico comando «Modifica».
    \item \textbf{Campo precompilato}: il campo di rinomina si apre con il nome
    corrente già inserito e selezionato, pronto per essere modificato o
    sostituito.
    \item \textbf{Rinomina dall'anteprima}: la possibilità di rinominare in linea
    direttamente dall'anteprima del documento, con un doppio click sul titolo.
  \end{itemize}
  In questo modo la riduzione del numero di click ottenuta con la
  riprogettazione, oggi vanificata dal costo della riscrittura manuale del nome,
  verrebbe finalmente sfruttata appieno.

  \item \textbf{Introduzione di funzionalità social.} Allo stato attuale
  l'interazione tra gli utenti è limitata ai soli commenti sui documenti
  condivisi nella sezione Community, senza alcuna possibilità di entrare in
  contatto diretto o di costruire una rete di relazioni; questa carenza
  indebolisce il senso di comunità e la collaborazione, aspetti particolarmente
  rilevanti in un contesto di studio condiviso. Si propone pertanto
  l'implementazione di funzionalità di tipo social, quali la possibilità di
  seguire e smettere di seguire altri utenti (follow/unfollow) e di scambiare
  messaggi diretti all'interno della piattaforma. Tali funzionalità
  rafforzerebbero la dimensione comunitaria già presente, favorendo la
  condivisione di materiale e conoscenze e incentivando un utilizzo più
  continuativo e collaborativo del sito.

  \item \textbf{Pianificazione intelligente dello studio.} Nella sezione Esami
  l'utente può collegare i propri materiali e definire manualmente un elenco di
  attività, ma è lasciato solo nell'organizzazione della preparazione: deve
  stabilire autonomamente quali materiali utilizzare, su quali argomenti
  concentrarsi e come distribuire lo studio nel tempo disponibile, senza alcun
  supporto che lo aiuti a pianificare in modo efficace. Si propone quindi
  l'integrazione di un sistema basato su intelligenza artificiale, articolato in
  due fasi:
  \begin{itemize}
    \item \textbf{Analisi integrata dei dati}: il sistema considera le date
    degli appelli, il carico di studio di ciascun esame e tutto il materiale
    disponibile, sia quello caricato dall'utente o presente sulla piattaforma,
    sia i dati storici e i commenti lasciati da altri utenti riguardo allo
    stesso esame.
    \item \textbf{Piano di studio personalizzato}: a partire da tale analisi, il
    sistema genera un piano che indica i materiali consigliati, gli argomenti su
    cui concentrarsi e un programma operativo giornaliero, calibrato sul tempo a
    disposizione, che ripartisce il carico in modo equilibrato e sostenibile.
  \end{itemize}
  Una funzionalità di questo tipo ridurrebbe lo sforzo di pianificazione
  richiesto allo studente, offrendo un punto di partenza concreto e organizzato
  per la preparazione e favorendo una gestione più consapevole ed efficiente del
  tempo di studio.

  \item \textbf{Monitoraggio dell'avanzamento dello studio.} Allo stato attuale
  la sezione Esami offre un'indicazione di progresso essenziale, legata al solo
  completamento delle singole attività, che non permette all'utente di percepire
  con chiarezza a che punto si trovi nella preparazione complessiva di un esame
  né di valutare se stia rispettando i tempi previsti. Si propone pertanto di
  arricchire la sezione con uno strumento di tracciamento visivo che, a partire
  dal piano generato dal sistema di intelligenza artificiale descritto in
  precedenza oppure da un piano definito manualmente, rappresenti graficamente
  sia la quota di preparazione completata sia l'aderenza ai tempi previsti. La
  possibilità di monitorare visivamente l'avanzamento renderebbe immediatamente
  percepibile il livello di preparazione raggiunto, aiutando l'utente a mantenere
  la motivazione, a individuare eventuali ritardi e a regolare di conseguenza il
  proprio ritmo di studio.

  \item \textbf{Sezione Workshop per lo studio assistito.} Le funzionalità
  attualmente offerte dalla piattaforma supportano la raccolta, l'organizzazione
  e la pianificazione del materiale, ma non l'attività di apprendimento in senso
  stretto: una volta predisposti gli appunti e definito il piano, l'utente è
  lasciato solo nello studio, privo di strumenti che lo aiutino a interrogare i
  contenuti, a individuarne i concetti chiave e a sintetizzarli. Si propone
  pertanto l'introduzione di una sezione dedicata, denominata \textit{Workshop},
  concepita come ambiente di studio attivo: selezionando un esame, tutto il
  materiale ad esso collegato verrebbe reso interrogabile ed elaborabile
  attraverso due strumenti complementari basati su intelligenza artificiale:
  \begin{itemize}
    \item \textbf{Interrogazione conversazionale dei contenuti}: l'utente può
    porre domande in linguaggio naturale e ricevere risposte fondate
    esclusivamente sul materiale dell'esame, riducendo il rischio di
    informazioni errate o non pertinenti.
    \item \textbf{Generazione automatica di materiali di sintesi}: a partire
    dallo stesso materiale, il sistema produce mappe concettuali, riassunti e
    schemi grafici a supporto della comprensione.
  \end{itemize}
  Un ambiente di questo tipo trasformerebbe la piattaforma da raccoglitore di
  materiali a strumento di apprendimento, favorendo una comprensione più rapida
  e approfondita dei contenuti e adattando il supporto alle effettive esigenze
  di studio dell'utente.

  \item \textbf{Sviluppo di un'app mobile dedicata.} Allo stato attuale la
  piattaforma è fruibile esclusivamente da web, una modalità adatta allo studio
  da postazione fissa ma che non sfrutta le potenzialità dei dispositivi mobili
  né accompagna lo studente nei contesti di mobilità. Si propone pertanto lo
  sviluppo di un'app mobile dedicata, che renda la piattaforma accessibile
  ovunque e ne arricchisca le funzionalità su due fronti:
  \begin{itemize}
    \item \textbf{Funzioni native in mobilità}: accesso immediato ai propri
    materiali, notifiche push relative a scadenze, piani di studio e avanzamento
    descritti in precedenza, e conversione automatica di registrazioni audio (ad
    esempio delle lezioni) in appunti testuali.
    \item \textbf{Condivisione e socialità}: caricamento e condivisione di
    documenti direttamente dal dispositivo e un'interazione più diretta e
    naturale con gli altri utenti, a complemento delle funzionalità social
    proposte in precedenza.
  \end{itemize}
  Un'app di questo tipo estenderebbe l'uso della piattaforma oltre la postazione
  fissa, integrando lo studio nella quotidianità dello studente e favorendo un
  utilizzo più continuativo e coinvolgente del servizio.
\end{itemize}
